\section{Modelos UML}

El conjunto completo de modelos reside en el ERS (\S\ref{sec:ers-final}) y como archivos individuales en \texttt{03\_Modelado/}, en formato editable y exportado. Esta secci\'on no los reproduce: los lee. La gu\'ia pide para este apartado una \emph{lectura interpretativa}, y eso es lo que sigue: qu\'e decisi\'on de ingenier\'ia sostiene cada modelo, qu\'e revela que los otros no muestran, y d\'onde los modelos se contradicen entre s\'i o con el texto.

\subsection{Qu\'e sostiene cada modelo}

La Tabla~\ref{tab:modelos-decision} resume, para cada modelo, la pregunta que responde y la decisi\'on de ingenier\'ia que sostiene.

\begin{table}[H]
\centering
\small
\begin{tabular}{|p{3.4cm}|p{5.6cm}|p{5.0cm}|}
\hline
\cellcolor{sigagreen}\textcolor{white}{\textbf{Modelo}} & \cellcolor{sigagreen}\textcolor{white}{\textbf{Pregunta que responde}} & \cellcolor{sigagreen}\textcolor{white}{\textbf{Decisi\'on que sostiene}} \\ \hline
Contexto & Qu\'e queda dentro y qu\'e fuera del sistema & Las cuatro exclusiones de RD-16 a RD-18 \\ \hline
Casos de uso (17) & Qu\'e puede hacer cada actor & El reparto de permisos de RF-19 \\ \hline
Clases refinadas (16) & Qu\'e entidades persisten y c\'omo se relacionan & La jerarqu\'ia \texttt{Equipo} y el acoplamiento que mide M5 \\ \hline
Secuencia (16) & En qu\'e orden colaboran los objetos & Los umbrales de latencia de NFR-01 y de los RF de control \\ \hline
Actividad (16) & Qu\'e decisiones y bifurcaciones tiene cada flujo & Los flujos alternativos de cada caso de uso \\ \hline
Estados (2) & Qu\'e ciclo de vida tienen las entidades no triviales & La columna \texttt{Estado} de la matriz de trazabilidad \\ \hline
Componentes & C\'omo se reparte la l\'ogica en m\'odulos & La estructura de m\'odulos del MVP \\ \hline
Despliegue & D\'onde se ejecuta cada parte & Las dependencias DEP-01 a DEP-09 \\ \hline
i* SD y SR & Por qu\'e cada actor quiere lo que quiere & La justificaci\'on de prioridad de los RF n\'ucleo \\ \hline
DFD nivel 1 & Por d\'onde circulan los datos & El eslab\'on \texttt{Proceso} de la cadena de trazabilidad \\ \hline
\end{tabular}
\caption{Modelos del sistema y decisi\'on de ingenier\'ia que cada uno sostiene}
\label{tab:modelos-decision}
\end{table}

\subsection{Lectura del modelo de clases}

El modelo refina diecis\'eis clases: \texttt{Aula}, \texttt{SensorIoT}, \texttt{LecturaSensor}, \texttt{Equipo}, \texttt{Proyector}, \texttt{Climatizador}, \texttt{GatewayIoT}, \texttt{Alerta}, \texttt{SolicitudMantenimiento}, \texttt{Usuario}, \texttt{Rol}, \texttt{HorarioAcademico}, \texttt{ConsumoEnergetico}, \texttt{ReporteAdministrativo}, \texttt{BitacoraAccion} y \texttt{CamaraVigilancia}.

La decisi\'on estructural m\'as relevante es la \textbf{jerarqu\'ia de herencia} \texttt{Equipo} $\rightarrow$ \{\texttt{Proyector}, \texttt{Climatizador}\}. No es cosm\'etica: es lo que permite que RF-02 (control remoto de dispositivos) se especifique una sola vez y que RF-04 y RF-05 sean casos concretos suyos sobre subtipos, en lugar de tres requisitos de control independientes que habr\'ia que mantener en paralelo. La alternativa descartada era modelar proyector y climatizador como clases sin relaci\'on, con un atributo \texttt{tipo} discriminante; se descart\'o porque obligaba a duplicar las operaciones de control y porque romp\'ia la extensibilidad que exige NFR-09 (incorporar un nuevo tipo de dispositivo sin tocar el n\'ucleo).

La segunda lectura relevante es de \textbf{acoplamiento}, y es la que explica la \'unica m\'etrica que no cumple. Dos clases concentran las relaciones: \texttt{Aula}, que participa en pr\'acticamente todo flujo de monitoreo, y \texttt{Equipo}, que participa en todo flujo de control. Por eso, al muestrear la modificabilidad (\S8), RF-01 y RF-07 arrastran seis requisitos cada uno. Ese acoplamiento no es un error de especificaci\'on sino una propiedad del dominio: un aula es, literalmente, el agregado que contiene todo lo dem\'as, pero tiene una consecuencia pr\'actica que conviene declarar: cualquier cambio en la definici\'on de \texttt{Aula} o en la de \texttt{LecturaSensor} obliga a revisar la mitad del cat\'alogo funcional.

\subsection{Lectura de los diagramas de estados}

Solo dos entidades tienen ciclo de vida modelado: \texttt{Alerta} (Generada $\rightarrow$ Notificada $\rightarrow$ \{Atendida $\mid$ Escalada\}) y \texttt{SolicitudMantenimiento} (Registrada $\rightarrow$ En atenci\'on $\rightarrow$ Cerrada). La pregunta pertinente no es por qu\'e hay dos, sino por qu\'e no hay m\'as, y la respuesta es un criterio expl\'icito: se modela el estado de una entidad cuando su comportamiento depende de su historia, no de su valor actual. Un \texttt{Aula} est\'a ocupada o vac\'ia seg\'un su \'ultima lectura, sin memoria; una \texttt{Alerta}, en cambio, no puede escalarse si no fue notificada antes. Aplicar el criterio al rev\'es (modelar un estado por cada clase) habr\'ia producido nueve diagramas triviales sin valor de verificaci\'on.

Esa decisi\'on tiene una consecuencia directa sobre la trazabilidad: la columna \texttt{Estado} de la matriz aparece como \emph{no aplica} en la mayor\'ia de las filas, y eso se declara con su raz\'on en cada una. La plantilla de la gu\'ia lo prev\'e expresamente al definir esa columna como aplicable ``si aplica''.

La transici\'on \texttt{Notificada} $\rightarrow$ \texttt{Escalada} merece atenci\'on aparte porque es la \'unica que se dispara por tiempo y no por acci\'on de un actor: ocurre cuando transcurre el umbral sin atenci\'on. Es, por tanto, el punto donde el modelo de estados y NFR-01 se tocan, y donde un cambio en el umbral de la alerta obliga a revisar ambos artefactos.

\subsection{Lectura del modelado organizacional i*}

Los diagramas i* responden una pregunta que ni los casos de uso ni las clases responden: \emph{por qu\'e} cada actor quiere lo que pide. El diagrama de Dependencia Estrat\'egica muestra que el Docente no depende del sistema, sino del Personal de Infraestructura, y que es este \'ultimo quien depende de SIGA para obtener datos de sensor en tiempo real. Esa cadena (Docente $\rightarrow$ Infraestructura $\rightarrow$ SIGA) explica un hecho que de otro modo parecer\'ia un descuido de elicitaci\'on: el Docente figura como origen de nueve requisitos pero casi nunca como usuario operativo del sistema. No es que se le haya ignorado; es que su dependencia est\'a mediada.

El diagrama de Raz\'on Estrat\'egica descompone el \emph{softgoal} ``reducir el tiempo de respuesta a fallas'' del Personal de Infraestructura en tres tareas que se corresponden con RF-07, RF-08 y RF-12. Esa descomposici\'on es la justificaci\'on documentada de por qu\'e esos tres requisitos son Must: no lo son por criterio del equipo, sino porque son las tres tareas en que se descompone el objetivo del actor con mayor poder e inter\'es del mapa de \emph{stakeholders}.

\subsection{CU-17: Ejercer derechos sobre datos personales (nuevo)}

RF-24 (exportaci\'on de datos personales) y RF-25 (rectificaci\'on) exist\'ian desde la Entrega 3, derivados del an\'alisis de los Art. 13 y 14 de la LOPDP, pero no ten\'ian caso de uso propio: colgaban impl\'icitamente de CU-11, cuyo prop\'osito es la gesti\'on de acceso por roles. Son cosas distintas: CU-11 responde a ``qui\'en puede entrar y a qu\'e'', mientras que el ejercicio de derechos responde a ``qu\'e puede exigir el titular sobre sus propios datos'', con actores, plazos legales y evidencia diferentes. Mientras esa confusi\'on persisti\'o, ambos requisitos figuraban como hu\'erfanos de caso de uso en el diagn\'ostico de trazabilidad.

CU-17 lo corrige: actor principal Docente, actores secundarios Personal Administrativo y de TI, con un curso normal que cubre las dos v\'ias (acceso y rectificaci\'on), el plazo de $\leq$15 d\'ias h\'abiles de NFR-11 y el registro en bit\'acora que permite demostrar cumplimiento ante la Autoridad de Protecci\'on de Datos. Su flujo alternativo A3 es deliberadamente conservador: cuando el plazo est\'a por vencer, el sistema alerta al responsable pero \textbf{no} resuelve la solicitud por su cuenta, porque una respuesta autom\'atica a un derecho del titular no ser\'ia defendible ante la autoridad.

\subsection{Diagrama de Flujo de Datos: Nivel 1 (nuevo)}

El ERS de la Entrega 3 modelaba estructura y comportamiento, pero ning\'un modelo describ\'ia el flujo de datos a nivel de proceso. Ese vac\'io no era cosm\'etico: el proceso DFD es un eslab\'on de la cadena de trazabilidad que la gu\'ia exige, y sin \'el ninguna cadena adelante pod\'ia cerrarse, que es la causa de que M4\_ade midiera 0\,\% en la primera auditor\'ia (\S8).

Los ocho procesos (Tabla~\ref{tab:dfd-procesos}) se derivaron agrupando los RF que comparten el mismo flujo de entrada y salida, verificado contra la matriz de trazabilidad; no provienen de una plantilla gen\'erica.

\begin{table}[H]
\centering
\small
\begin{tabular}{|l|p{5.5cm}|p{5.2cm}|}
\hline
\cellcolor{sigagreen}\textcolor{white}{\textbf{Proceso}} & \cellcolor{sigagreen}\textcolor{white}{\textbf{Nombre}} & \cellcolor{sigagreen}\textcolor{white}{\textbf{RF que agrupa}} \\ \hline
P-1 & Captura de datos ambientales y ocupaci\'on & RF-01, RF-03, RF-20, RF-22 \\ \hline
P-2 & Control remoto de equipos & RF-02, RF-04, RF-05 \\ \hline
P-3 & Generaci\'on y gesti\'on de alertas & RF-08, RF-11, RF-21 \\ \hline
P-4 & Automatizaci\'on energ\'etica & RF-13, RF-15, RF-16 \\ \hline
P-5 & Gesti\'on de mantenimiento & RF-10, RF-12 \\ \hline
P-6 & Reportes administrativos & RF-14, RF-17, RF-18 \\ \hline
P-7 & Gesti\'on de acceso y auditor\'ia & RF-19, RF-23 \\ \hline
P-8 & Gesti\'on de derechos sobre datos personales & RF-24, RF-25 \\ \hline
\end{tabular}
\caption{Procesos del DFD de nivel 1}
\label{tab:dfd-procesos}
\end{table}

El diagrama hace visibles dos cosas que ning\'un modelo anterior mostraba. La primera es que el almac\'en de bit\'acora (D4) recibe escrituras de \textbf{seis de los ocho procesos}. Es coherente con RF-23 y con el deber de responsabilidad demostrada de la LOPDP, pero convierte a ese almac\'en en un punto \'unico de fallo para la auditabilidad de todo el sistema: si la escritura en bit\'acora falla silenciosamente, el sistema sigue operando con normalidad y la p\'erdida solo se descubre cuando alguien pide la evidencia. Ning\'un requisito actual cubre ese escenario, y se deja anotado como candidato para la siguiente iteraci\'on.

La segunda es que P-8 no exist\'ia como proceso en ninguna vista previa, lo que confirma desde un \'angulo independiente el mismo problema que motiv\'o CU-17: el ejercicio de derechos del titular estaba disuelto dentro de la gesti\'on de acceso.

\subsection{Coherencia entre modelos}

Un conjunto de modelos puede ser individualmente correcto y colectivamente incoherente, de modo que se verificaron tres correspondencias que deb\'ian cumplirse por construcci\'on:

\begin{itemize}
\item \textbf{Casos de uso frente a clases.} Toda clase del modelo de dominio participa en al menos un caso de uso, y todo caso de uso referencia al menos una clase. Se comprob\'o sobre la matriz de trazabilidad, columna \texttt{Clase\_UML}.
\item \textbf{Casos de uso frente a actores.} Los nueve actores del diagrama general participan en al menos un caso de uso, desde Personal de Infraestructura, presente en diez, hasta el Sistema de Videovigilancia, presente solo en CU-15. Ese conteo es el que sustenta la m\'etrica M1c (\S8).
\item \textbf{Procesos DFD frente a requisitos.} Los 25 RF est\'an cubiertos por exactamente uno de los ocho procesos, sin solapamiento ni hueco. Los RF de IA se adscriben al proceso del RF que refinan.
\end{itemize}

La \'unica incoherencia detectada entre modelos y texto durante esta verificaci\'on fue la ausencia f\'isica de una figura de diagrama de estados que el ERS referenciaba (defecto D-11): el archivo exist\'ia en la carpeta de modelado pero no en la de figuras del documento, de modo que el modelo estaba, y lo que faltaba era su v\'inculo.
